\section{Introduction}
% DRAFT 2 (2026-09-27). Budget ~1 page. Funnel: DH today -> OCR limits -> damage hits facts -> GenAI temptation
% and risk -> gap -> RQ -> this experiment -> contributions. FACT SOURCE: notes/report_sources.md Sec. 1-4.
% Every number and quote below was re-checked against the source PDF in ref_dump/ on 2026-09-27 (not just the
% abstract). Two citations replaced this pass: bingham2010digitization and holley2009good were dropped from the
% Zotero library (not independently accessible), so their claims now use nguyen2021survey and booth-etal-2024-bln600,
% both already cited elsewhere in this section and both open-access PDFs in ref_dump/.

Digitised newspapers are now a core source for historical research. Digitisation lets scholars search and mine
large collections by keyword rather than reading page by page \citep{nguyen2021survey}. Yet what historians search
and mine is not the page itself but the text that an OCR engine read from a scan of it. OCR errors measurably change
the results of quantitative analyses \citep{hill2019quantifying} and harm downstream NLP tasks \citep{vanstrien2020assessing},
with named entities among the most affected \citep{hamdi2023indepth}. For 19th-century British newspapers the mean
character error rate (CER) is about 7\%, but it ranges from 0.3\% to 44.5\% between excerpts of a single title, and
this spread tracks the condition of the source page, not just the OCR engine: faded or physically damaged originals
give the engine the least to work with \citep{booth-etal-2024-bln600}. For the British Library's own 19th-century
newspaper archive, character accuracy reaches 83.6\% but word accuracy falls to 78\%, and accuracy for capitalised
words and number groups drops further still, to 63.4\% and 64.1\% \citep{tanner2009measuring}. The words hit hardest
are thus often the facts a historian searches for: names, places, dates and sums of money.

Generative AI makes it tempting to repair such text automatically, but the evidence is mixed. Llama~2,
instruction-tuned on BLN600, reduced the CER of its OCR by 54.5\%, against 23.3\% for a fine-tuned BART
\citep{thomas2024leveraging}, while a study of 14 LLMs on several historical benchmarks found them ``anything but
efficient'' at post-correction \citep{boros2024postcorrection}. There is also a deeper risk. A generative model always
returns fluent text, even where the evidence for a word is gone, and such models are known to produce content that
their input does not support \citep{ji2023survey}. For an archive, a plausible but wrong name is worse than visible
noise, because the reader can no longer see that anything was damaged. Similarity-based metrics may hide this error
too: a sentence with the wrong name can still score as close in meaning to the original.

Existing post-OCR correction studies score whole documents at whatever error level the OCR happened to produce. They show that a
model helps on average, but not how far repair can be trusted as damage grows, nor whether the facts inside the damaged
region survive. We therefore ask:

\begin{quote}
\textbf{RQ.} When the damaged words of a sentence are known and their garbled letters are visible, how much of the
original wording, and in particular its names, numbers and dates, can a fine-tuned masked language model and a
few-shot instruction-following LLM restore, and how does this change as the damaged share of the sentence grows?
\end{quote}

\noindent We split it into three sub-questions. \textbf{SQ1} (severity): does repair quality fall smoothly as damage
grows, or break down at some level? \textbf{SQ2} (method): given identical input, does a fine-tuned masked language
model or a few-shot LLM repair better? \textbf{SQ3} (metric validity): does semantic similarity overstate repair
quality compared with exact fact recovery?

Answering these questions needs damage whose amount we control and whose original text we know, which real damaged
pages cannot give. We therefore build a controlled experiment on BLN600 \citep{booth-etal-2024-bln600}, a corpus of
600 newspaper excerpts with human gold transcriptions. In each of 150 test sentences we damage one block of words
around a fact at five nested levels, from one word to 75\% of the sentence, with character noise calibrated on the
corpus's own OCR errors. Both methods receive the same input: the intact context plus the garbled form of each damaged
word, so that only the repair differs. We compare two GenAI families: masked infilling with BERT
\citep{devlin2019bert}, fine-tuned on in-domain gold text and guided by letter similarity in the spirit of
``context beats confusion'' \citep{evershed2014correcting}, and instruction-following generation with an open-weight
LLM prompted with three solved examples. We score the repaired words with BERTScore \citep{zhang2020bertscore} for
meaning and with a Fact Recovery Rate for exact facts, and we probe the LLM for memorised test text
\citep{sainz2023contamination}. Repair here works on text only: we simulate the OCR output that damage produces, not
the physical page, and we assume the damaged region is known; detecting it, and repairing directly from the scanned
image with a vision-capable model, are both left to future work.

Our contributions are (i) a fact-centred, graded damage benchmark on BLN600 with nested spans, so the same place
gets worse from level to level; (ii) a controlled comparison of masked infilling and few-shot generation under
identical input and scoring; and (iii) a fact-level metric next to BERTScore that shows where semantic similarity and
factual accuracy part ways. \todo{Key outcomes, from the TEST split only: (1) the level at which Fact Recovery Rate
falls to the no-repair floor; (2) which method wins at low vs high damage, and whether a crossing is a format-failure
effect; (3) the size of the BERTScore vs Fact Recovery gap.}

% ---- OPEN POINTS ----
%  1. RESOLVED 2026-09-27: "IS BERT GENAI?" answered in Related Work (Sec. 2, citing Wang and Cho 2019 for
%     "BERT can generate", plus the fallback reading of BERT as the in-domain baseline the LLM must beat).
%  2. KNOWN LOCATION is a strong assumption; justified in Method (OCR engines flag low-confidence regions).
%  3. RESOLVED 2026-09-27: two \todo{get citation} gaps closed. Content-checked against full PDF text
%     (not just abstract) in ref_dump/: nguyen2021survey's own introduction states digitisation is "for the
%     purpose of text processing by computers (e.g., search...), ... easier access ... by a much wider
%     audience"; booth-etal-2024-bln600 Sec 3.2 states BLN600 source images include some "too faded or
%     damaged to read" and that "older newspapers suffer with quality issues ... particularly those with
%     faded low-contrast print." Also fixed: vanstrien2020assessing was cited for "named entities among
%     the most affected", but that paper actually finds OCR errors have a SMALLER impact on NER than on
%     sentence segmentation (its own words). Decoupled: vanstrien2020assessing now backs only the general
%     "harm downstream NLP tasks" claim (accurate: that is the paper's own subject), and hamdi2023indepth
%     alone backs "named entities among the most affected" (verified: "NEs are particularly affected").
